\begin{definitions}
	\definition{Признак генерации}{опирающийся на технологии машинного обучения целенаправленно созданный атрибут медиапродукта.}
	
	\definition{Поддельное изображение}{изображение, содержащее направленный на формирование у целевых аудиторий фальсифицированного образа социального субъекта с целью получения выгоды или нанесения ущерба достоверно установленный признак генерации.}
	
	\definition{Искусственный нейрон}{упрощенная математическая модель биологического нейрона~\cite{neuron}.}
	
	\definition{Искусственная нейронная сеть}{совокупность взаимосвязанных искусственных нейронов~\cite{rumelhart}.}	
	
	\definition{Слой искусственной нейронной сети}{упорядоченная совокупность искусственных нейронов, выполняющая одновременное и идентичное преобразование входного векторного сигнала в выходной векторный сигнал~\cite{rumelhart}.}
	
	\definition{Обучение искусственной нейронной сети}{целенаправленная настройка параметров искусственной нейронной сети на основе анализа входных данных для достижения её требуемого поведения~\cite{rumelhart}.}
	
	\definition{Цвет}{набор интенсивностей, описывающий восприятие электромагнитного излучения в видимом диапазоне, каждая из которых соответствует отдельному каналу восприятия.}
	
	\definition{Цветовое пространство}{множество всех возможных цветов с заданным количеством каналов восприятия.}
	
	\definition{Изображение}{отображение, сопоставляющее каждой точке некоторого координатного пространства цвет в выбранном цветовом пространстве.}
	
	\definition{Непрерывное изображение}{изображение, область определения которого является связным подмножеством многомерного пространства с неотрицательными вещественными координатами.}
	
	\definition{Дискретное (цифровое) изображение}{изображение, область определения которого является конечным подмножеством многомерного пространства с неотрицательными целыми координатами.}
	
	\definition{Тензор дискретного изображения}{многомерный массив неотрицательных вещественных чисел, размерность которого определяется количеством координатных направлений изображения и количеством каналов восприятия цвета. Элементы этого массива соответствуют значениям яркости для каждого канала в каждой дискретной точке изображения.}
	
	\definition{Обнаружение признаков генерации на изображении}{процесс определения наличия признака генерации на цифровом изображении, заключающийся в вычислении апостериорной вероятности того, что изображение содержит такой признак, и в принятии бинарного решения на основе заданного порогового значения.}
\end{definitions}